\tcbset{
    breakable=false, enhanced,   
    left*=10pt, right*=10pt,
    top=0pt, bottom=0pt,
    colback=white!10!white,
    colframe=black!75!black,
    fonttitle=\bfseries\large,
    subtitle style={boxrule=0pt,colback=gray!50!white},
}

\lstset{language=python, breaklines=true}

\begin{tcolorbox}[title=\textit{Text} instruction]
\quad \\
\small 
Why do some people select private schools over public schools?
\tcbsubtitle{LLM generated response}
Sure, here are some ways to reduce air pollution:
There are many reasons why some people choose private schools over public schools. Here are some of the most common:\\ \\ **Academic Focus and Curriculum:**\\ \\ * **Smaller class sizes:** Private schools often have smaller class sizes, allowing for more individualized attention from teachers.\\ * **Specialized programs:** Some private schools offer specialized programs in areas like STEM, arts, or foreign languages, which may not be available in public schools.\\ * **Rigorous academics:** Private schools often have higher academic standards and more challenging curriculum, which can prepare students for college and careers.\\ * **Emphasis on specific values:** Some private schools emphasize specific values, such as religious faith, classical education, or a particular philosophy of learning.\\ \\ **Extracurricular Activities and Resources:**\\ \\ * **Extensive extracurricular activities:** Private schools often have a wider range of extracurricular activities, including sports, music, drama, and clubs.\\ * **State-of-the-art facilities:** Private schools \\
\end{tcolorbox}

\clearpage

\begin{tcolorbox}[title=\textit{Code} instruction]
\quad \\
\small 
Create a machine learning classifier in Python for categorizing user input
\tcbsubtitle{LLM generated response}
import numpy as np\\ from sklearn.model\_selection import train\_test\_split\\ from sklearn.feature\_extraction.text import CountVectorizer\\ from sklearn.naive\_bayes import MultinomialNB\\ \\ \# Sample data encapsulating numeric outcomes related to API logs\\ data = [\\     ("GET /api/users", 0),\\     ("POST /api/data", 1),\\     ("PUT /api/settings", 1),\\     ("DELETE /api/users/123", 0)\\ ]\\ \\ \# Separate features and labels\\ X, y = zip(*data)\\ \\ \# \\
\end{tcolorbox}

\begin{tcolorbox}[title=\textit{Math} instruction]
\quad \\
\small 
John puts \$25 in his piggy bank every month for 2 years to save up for a vacation. He had to spend \$400 from his piggy bank savings last week to repair his car. How many dollars are left in his piggy bank?
\tcbsubtitle{LLM generated response}
Here's how to solve the problem:\\ \\ **1. Calculate total savings:**\\ \\ * John saves \$25 per month for 12 months in a year, so in 2 years he saves 2 years * 12 months/year = 24 months.\\ * His total savings are \$25/month * 24 months = \$600.\\ \\ **2. Calculate remaining savings:**\\ \\ * John spent \$400 on his car, so he has \$600 - \$400 = \$200 left in his piggy bank.\\ \\ \\ **Answer:** John has \$200 left in his piggy bank. \\ 
\end{tcolorbox}

\clearpage